
\documentclass[12pt]{article}
\usepackage{amsmath,amssymb}
\usepackage{fullpage}
\usepackage{listings}
\usepackage{enumitem}
\usepackage{xcolor}
\usepackage{courier}

\lstset{
    basicstyle=\ttfamily\small,
    breaklines=true,
    showstringspaces=false,
    columns=fullflexible,
    frame=single,
    xleftmargin=0.1in,
    xrightmargin=0.1in
}

\title{CPSC 250 -- Programming for Data Manipulation\\Final Exam Solutions}
\date{}

\begin{document}

\maketitle

\section*{Part I: Multiple Choice}

\begin{enumerate}[label=\arabic*.]
\item C
\item B
\item C
\item C
\item A
\item B
\item A
\item B
\item B
\item C
\item A
\item B
\item B
\item A
\item B
\end{enumerate}

\section*{Part II: Error Identification and Correction}

One corrected version is shown below.

\begin{lstlisting}[language=Python]
def above_average(values):
    total = 0
    for value in values:
        total += value

    average = total / len(values)

    result = []
    for value in values:
        if value > average:
            result.append(value)

    return result

measurements = [4.2, 5.1, 3.8, 6.0, 4.9]
print(above_average(measurements))
\end{lstlisting}

Important errors include:

\begin{itemize}
\item The loop should accumulate the values using \texttt{total += values[i]} or \texttt{total += value}; the original line overwrites \texttt{total} each time.
\item The \texttt{if} statement needs a colon.
\item \texttt{append} is a method and should be called with parentheses, not brackets.
\item The variable name in the final line should be \texttt{measurements}, not \texttt{measurement}.
\end{itemize}

\section*{Part III: Code Writing}

\subsection*{Q1. Functions and collections of objects}

One possible solution is:

\begin{lstlisting}[language=Python]
def average_height_with_sunlight(plants, min_sunlight):
    total = 0
    count = 0

    for plant in plants:
        if plant.sunlight >= min_sunlight:
            total += plant.height
            count += 1

    if count == 0:
        return 0

    return total / count
\end{lstlisting}

A solution using a temporary list is also fine:

\begin{lstlisting}[language=Python]
def average_height_with_sunlight(plants, min_sunlight):
    heights = []

    for plant in plants:
        if plant.sunlight >= min_sunlight:
            heights.append(plant.height)

    if len(heights) == 0:
        return 0

    return sum(heights) / len(heights)
\end{lstlisting}

\subsection*{Q2. Classes and collections of objects}

\begin{lstlisting}[language=Python]
class Book:

    def __init__(self, title, author, pages):
        self.title = title
        self.author = author
        self.pages = pages

    def __str__(self):
        return f"{self.title} by {self.author} ({self.pages} pages)"

    def is_long(self):
        return self.pages > 300


for book in library:
    if book.is_long():
        print(book)
\end{lstlisting}

\subsection*{Q3. Inheritance and polymorphism}

\begin{lstlisting}[language=Python]
class Account:

    def __init__(self, owner, balance):
        self.owner = owner
        self.balance = balance

    def monthly_update(self):
        return self.balance


class SavingsAccount(Account):

    def __init__(self, owner, balance, rate):
        super().__init__(owner, balance)
        self.rate = rate

    def monthly_update(self):
        self.balance = self.balance * (1 + self.rate)
        return self.balance


accounts = [
    Account("Alice", 1000),
    SavingsAccount("Bob", 1000, 0.01)
]

for account in accounts:
    print(account.monthly_update())
\end{lstlisting}

The loop demonstrates polymorphism because it calls \texttt{monthly\_update()} on each object without needing to know whether the object is an \texttt{Account} or a \texttt{SavingsAccount}.

\subsection*{Q4. pandas, plotting, and model functions}

\begin{lstlisting}[language=Python]
import pandas as pd
import matplotlib.pyplot as plt

df = pd.read_csv("experiment.csv")

valid = df[df["valid"] == True]

plt.scatter(valid["time"], valid["position"])
plt.xlabel("time")
plt.ylabel("position")

def linear(x, m, b):
    return m*x + b
\end{lstlisting}

It would also be acceptable to write the filter as:

\begin{lstlisting}[language=Python]
valid = df[df["valid"]]
\end{lstlisting}

\section*{Part IV: Code Comprehension and Commenting}

Possible comments are shown below. Student answers do not need to match this wording exactly, but they should explain the purpose of each section or function.

\begin{lstlisting}[language=Python]
# Import the libraries needed for data analysis and plotting.

import pandas as pd
import matplotlib.pyplot as plt

# Define a class to represent a plant in the data set.

class Plant:

    # Initialize a new Plant object with its attributes.
    def __init__(self, name, height, sunlight):
        self.name = name
        self.height = height
        self.sunlight = sunlight

    # Return a readable string description of the plant.
    def __str__(self):
        return f"{self.name}: {self.height} cm, {self.sunlight} hours"

    # Determine whether the plant needs additional attention.
    def needs_attention(self):
        return self.height < 10 or self.sunlight < 4


# Read the CSV file and create a list of Plant objects.

plants = []

df = pd.read_csv("plants.csv")

for _, row in df.iterrows():
    plants.append(
        Plant(row["Name"], row["Height"], row["Sunlight"])
    )


# Select and print the plants that need attention.

attention = [p for p in plants if p.needs_attention()]

for p in attention:
    print(p)


# Create a DataFrame and make a bar chart of plant heights.

plot_df = pd.DataFrame({
    "Name": [p.name for p in plants],
    "Height": [p.height for p in plants]
})

plt.bar(plot_df["Name"], plot_df["Height"])
plt.xlabel("Plant")
plt.ylabel("Height (cm)")
plt.title("Plant Heights")
plt.show()
\end{lstlisting}

\subsection*{Questions}

\begin{enumerate}[label=\arabic*.]
\item After the \texttt{for} loop completes, the \texttt{plants} list stores \texttt{Plant} objects.

\item The program is an example of using a collection of objects because it stores many \texttt{Plant} objects in a list and then processes those objects using loops and list comprehensions.

\item A plant is added to the \texttt{attention} list when \texttt{p.needs\_attention()} returns \texttt{True}. In this program, that means the plant's height is less than 10 or its sunlight value is less than 4.

\item The list comprehensions near the end extract the \texttt{name} and \texttt{height} attributes from each \texttt{Plant} object so that those values can be placed into a pandas DataFrame for plotting.

\item The main loop could still work because of polymorphism. If a \texttt{Flower} object inherits from \texttt{Plant} and provides its own version of \texttt{needs\_attention()}, then the code can still call \texttt{p.needs\_attention()} without needing major changes.
\end{enumerate}

\section*{Part V: Short Answer / Conceptual Design}

\begin{enumerate}[label=\arabic*.]
\item Choosing the right representation is important because the representation determines what operations are natural, efficient, and easy to express. For example, a raw CSV file is useful for storage, a pandas DataFrame is useful for filtering and aggregating data, a list of objects is useful when each item has related attributes and behaviors, and a plot or model can summarize patterns. Good data representation makes the algorithm simpler and makes the program easier to understand, test, and extend.

\item Refactoring into functions and classes makes a program easier to test because each function or method can be checked separately. It also makes debugging easier because the source of an error is usually confined to a smaller piece of code. Classes help organize related data and behavior, while functions reduce repeated code. This makes the program easier to extend because a change can often be made in one place rather than copied into many repeated blocks.
\end{enumerate}

\end{document}
